\documentclass[a4paper,10pt]{article}
\usepackage[T1]{fontenc}
\usepackage[utf8]{inputenc}
\usepackage[french]{babel}
\usepackage{lmodern}
\usepackage[margin=1.8cm]{geometry}
\usepackage{microtype}
\usepackage{enumitem}
\usepackage{titlesec}
\usepackage[hidelinks]{hyperref}
\setlength{\parindent}{0pt}
\setlength{\parskip}{4pt}
\titlespacing*{\section}{0pt}{9pt}{4pt}
\titleformat{\section}{\normalsize\bfseries}{\thesection.}{.5em}{}
\setlist[itemize]{leftmargin=*,nosep,topsep=3pt}
\pagestyle{plain}

\begin{document}
\begin{center}
  {\Large\bfseries État de l'art des attaques contre les ML-NIDS}\par
  \vspace{3pt}
\end{center}
\vspace{-5pt}

\section{Périmètre et classification}
Un système de détection d'intrusions réseau fondé sur l'apprentissage automatique
(\emph{ML-NIDS}) transforme des paquets en caractéristiques : durées, tailles,
nombres de paquets, indicateurs TCP ou relations entre machines. Un modèle
(SVM, réseau de neurones, autoencodeur ou réseau sur graphes) utilise cette
représentation pour reconnaître les activités malveillantes. Les attaques
adversariales visent cette chaîne de détection : elles doivent être distinguées
de l'activité réseau sous-jacente, par exemple un scan, un déni de service ou
une force brute, qu'elles cherchent à dissimuler.

Deux grandes familles se distinguent par leur moment d'action :
\begin{itemize}
  \item \textbf{Évasion :} à l'inférence, modifier une communication malveillante
  pour qu'un détecteur déjà entraîné la considère comme bénigne.
  \item \textbf{Empoisonnement :} altérer les données ou leurs étiquettes avant
  un entraînement ou un réentraînement pour altérer les décisions futures.
  Les \textbf{portes dérobées} constituent une forme ciblée d'empoisonnement
  qui apprend au
  modèle à produire une décision choisie lorsqu'un \emph{déclencheur} apparaît,
  tout en conservant de bonnes performances sur les données sans déclencheur.
\end{itemize}
Cette classification se croise avec la connaissance de l'adversaire :
\textbf{boîte blanche} (accès au modèle et à ses paramètres), \textbf{boîte grise}
(informations partielles) ou \textbf{boîte noire} (sans accès interne).
En boîte noire, l'adversaire peut exploiter des réponses du détecteur ou
transférer des perturbations préparées sur un modèle substitut
\cite{han,mao,zhang}. Il faut également préciser son contrôle : trafic émis,
machines compromises, données d'apprentissage ou accès aux prédictions.

\section{Évasion : des caractéristiques aux communications}
\textbf{Perturbation de caractéristiques.} TIKI-TAKA compare cinq attaques en
boîte noire (NES, Boundary, Pointwise, HopSkipJump et OPT) contre des MLP, CNN
et C-LSTM \cite{zhang}. Les requêtes guident la recherche d'un vecteur classé
comme bénin ; une variante restreint les modifications aux variables
temporelles. Mao et al. proposent SPTS, qui impose des contraintes et des
relations entre caractéristiques, puis DGM, qui améliore le transfert entre
modèles \cite{mao}. Leur démarche inclut un passage aux paquets. Toutefois,
un vecteur adversarial valide mathématiquement ne garantit pas qu'une
communication correspondante soit réalisable et conserve son effet malveillant.

\textbf{Remodelage du trafic.} TrafficManipulator agit sur les délais et ajoute
des paquets pour modifier les statistiques observées, notamment contre
KitNET, un ensemble d'autoencodeurs \cite{han}. TANTRA apprend les intervalles
entre paquets du trafic bénin avec un LSTM, puis adapte la temporisation du
trafic malveillant sans en changer le contenu \cite{tantra}. Ce LSTM est le
générateur de l'attaque ; les détecteurs ciblés comprennent KitNET, un
autoencodeur et Isolation Forest. Le coût en latence reste une contrainte,
en particulier pour les activités dépendantes du débit.

\textbf{Transformation guidée par l'explicabilité.} Okada et al. utilisent
KernelSHAP pour identifier les caractéristiques influentes d'un MLP en boîte
noire, puis les modifier par des actions sur le trafic \cite{okada}.
Selon l'extracteur, une force brute FTP est adaptée en modifiant les indicateurs
PSH ou en séparant les tentatives dans des sessions TCP distinctes. L'étude
évalue aussi des attaques XSS entre machines virtuelles et vérifie la
préservation de l'activité malveillante : elle relie explicitement évasion du
modèle et fonctionnement de l'attaque réseau.

\textbf{Perturbation structurelle.} Pour E-GraphSAGE et LineGraphSAGE, Venturi
et al. modifient les communications représentées par les arêtes et les nœuds
du graphe \cite{venturi}. La variante C2x\_B ajoute des communications bénignes
depuis un hôte compromis, sans retirer le flux malveillant : le voisinage
observé change et peut tromper le modèle. Ces expériences sur flux et graphes
ne constituent cependant pas une validation complète en réseau actif.

\newpage
\section{Empoisonnement et portes dérobées}
\textbf{Dégradation de l'apprentissage.} Venkatesan et al. étudient un SVM
distinguant trafic bénin et scans Nmap \cite{venkatesan}. L'adversaire insère
des exemples étiquetés « scan » près d'exemples bénins reconnus avec une forte
confiance. Lors du réentraînement, ces contradictions déplacent la frontière
de décision et augmentent les erreurs, notamment les faux positifs.
L'attaque vise ici la disponibilité du détecteur, plutôt que seulement
l'évasion d'un flux particulier. L'étude impose des contraintes de cohérence
aux caractéristiques, mais ne démontre pas une chaîne complète de rejeu réseau.
Une condition essentielle est donc l'accès à une partie des données
d'apprentissage : envoyer du trafic à un modèle figé ne suffit pas à
l'empoisonner.

\textbf{Comportement conditionné par un déclencheur.} PCAP-Backdoor insère des
paquets construits dans une partie du trafic bénin utilisé pour l'entraînement,
en conservant les étiquettes initiales : c'est un empoisonnement à
\emph{étiquettes propres} \cite{backdoor}. Le modèle apprend une association
entre les statistiques induites par ce déclencheur et la classe bénigne.
À l'inférence, l'ajout du déclencheur au trafic malveillant provoque alors une
classification bénigne. À la différence de l'empoisonnement précédent,
les résultats sur trafic sans déclencheur peuvent rester bons : une évaluation
standard ne suffit donc pas à révéler la compromission. L'étude agit sur les
PCAP et examine plusieurs modèles et jeux de données IoT et industriels.

\section{Enseignements, défenses et critères d'évaluation}
Le corpus montre que la vulnérabilité dépend du \textbf{couple modèle--extracteur}
et des capacités de l'adversaire. Les GNN introduisent une surface d'attaque
structurelle ; les détecteurs fondés sur des statistiques temporelles peuvent
être sensibles à des changements de rythme. Les résultats obtenus sur Kitsune,
CICIDS, TON\_IoT ou CTU-13 restent liés au protocole expérimental : des taux de
succès issus d'études différentes ne sont pas directement comparables.

Les défenses étudiées comprennent l'entraînement adversarial et les ensembles
de modèles, la détection des requêtes répétées \cite{han,zhang}, ainsi que la
sanitisation des données par désaccord entre classifieurs \cite{venkatesan}.
PCAP-Backdoor souligne également les limites du regroupement des activations
pour identifier certains déclencheurs \cite{backdoor}. Ces méthodes atténuent
des scénarios précis, sans établir une robustesse universelle.

Pour le démonstrateur, une expérience doit comparer un état de référence et un
état attaqué, avec le \textbf{rappel du trafic malveillant}, les \textbf{faux
positifs} et le \textbf{coût de l'attaque} : délai, volume, nombre de requêtes ou
fraction de données empoisonnées. Le taux d'évasion doit préciser son
dénominateur, notamment si seuls les flux initialement détectés sont retenus.
Pour une porte dérobée, il faut mesurer séparément le succès avec déclencheur
et la qualité sans déclencheur. Enfin, les caractéristiques doivent être
réextraites après transformation des paquets et la réussite de l'activité réseau
doit être vérifiée. TrafficManipulator/KitNET, TANTRA et les scénarios d'Okada
constituent ainsi des pistes complémentaires pour illustrer l'évasion ; le SVM
empoisonné et PCAP-Backdoor couvrent les risques liés à l'apprentissage.

\begingroup
\small
\setlength{\parskip}{0pt}
\renewcommand{\refname}{Références du corpus}
\begin{thebibliography}{8}
\setlength{\itemsep}{1pt}
\bibitem{han} D. Han et al. (2021). \emph{Evaluating and Improving Adversarial
Robustness of Machine Learning-Based Network Intrusion Detectors}. IEEE JSAC.
\href{https://doi.org/10.1109/JSAC.2021.3087242}{doi:10.1109/JSAC.2021.3087242}.
\bibitem{tantra} Y. Sharon et al. (2022). \emph{TANTRA: Timing-Based Adversarial
Network Traffic Reshaping Attack}. IEEE TIFS.
\href{https://doi.org/10.1109/TIFS.2022.3201377}{doi:10.1109/TIFS.2022.3201377}.
\bibitem{okada} S. Okada et al. (2025). \emph{Xai-driven black-box adversarial
attacks on network intrusion detectors}. IJIS.
\href{https://doi.org/10.1007/s10207-025-01016-0}{doi:10.1007/s10207-025-01016-0}.
\bibitem{venturi} A. Venturi, D. Stabili et M. Marchetti (2024).
\emph{Problem space structural adversarial attacks for Network Intrusion
Detection Systems based on Graph Neural Networks}.
\href{https://arxiv.org/abs/2403.11830}{arXiv:2403.11830v2}.
\bibitem{mao} Z. Mao et al. (2025). \emph{Exploring transferable adversarial
attacks for Deep Learning-based Network Intrusion Detection}. JNCA.
\href{https://doi.org/10.1016/j.jnca.2025.104255}{doi:10.1016/j.jnca.2025.104255}.
\bibitem{zhang} C. Zhang, X. Costa-Pérez et P. Patras (2022).
\emph{Adversarial Attacks Against Deep Learning-Based Network Intrusion
Detection Systems and Defense Mechanisms}. IEEE/ACM ToN.
\href{https://doi.org/10.1109/TNET.2021.3137084}{doi:10.1109/TNET.2021.3137084}.
\bibitem{venkatesan} S. Venkatesan et al. (2021). \emph{Poisoning Attacks and
Data Sanitization Mitigations for Machine Learning Models in Network Intrusion
Detection Systems}. MILCOM.
\href{https://doi.org/10.1109/MILCOM52596.2021.9652916}{doi:10.1109/MILCOM52596.2021.9652916}.
\bibitem{backdoor} A. Koyatan Chathoth et al. (2026). \emph{PCAP-Backdoor:
Backdoor Generator in Network Traffic for Intrusion Detection Systems}.
ACM TCPS. \href{https://doi.org/10.1145/3805033}{doi:10.1145/3805033}.
\end{thebibliography}
\endgroup
\end{document}
